\documentclass[11pt,a4paper]{article}

% ---------------------------------------------------------------------------
% Packages
% ---------------------------------------------------------------------------
\usepackage{amsmath,amssymb}
\usepackage{physics}
\usepackage{braket}
\usepackage[colorlinks=true,linkcolor=blue,citecolor=blue,urlcolor=blue]{hyperref}
\usepackage{graphicx}
\usepackage{booktabs}
\usepackage{enumitem}
\usepackage[margin=2.5cm]{geometry}
\usepackage{fancyhdr}

% ---------------------------------------------------------------------------
% Page style
% ---------------------------------------------------------------------------
\pagestyle{fancy}
\fancyhf{}
\fancyhead[L]{Pion Soft Factor -- Physics Note}
\fancyhead[R]{\today}
\fancyfoot[C]{\thepage}
\renewcommand{\headrulewidth}{0.4pt}

% ---------------------------------------------------------------------------
% Math convenience
% ---------------------------------------------------------------------------
\newcommand{\Gfive}{\gamma_{5}}
\newcommand{\Gmu}{\gamma_{\mu}}
\newcommand{\Gnu}{\gamma_{\nu}}
\newcommand{\Grho}{\gamma_{\rho}}
\newcommand{\Gsigma}{\gamma_{\sigma}}
\newcommand{\Gone}{\gamma_{1}}
\newcommand{\Gtwo}{\gamma_{2}}
\newcommand{\Gthree}{\gamma_{3}}
\newcommand{\einsum}{\texttt{einsum}}
\newcommand{\xp}{\texttt{xp}}
\newcommand{\conj}[1]{#1^{\dagger}}
\newcommand{\gbar}[1]{\bar{#1}}
\newcommand{\walltag}[1]{\texttt{#1}}

% Dirac matrix shorthand in text
\newcommand{\gI}{\mathbf{1}}

\begin{document}

% ===========================================================================
\title{\textbf{Pion Soft Factor Measurement}\\[4pt]
\large A Detailed Physics Note for the PyQUDA Implementation}
\author{}
\date{\today}
\maketitle
\thispagestyle{fancy}

% ===========================================================================
\section{Physics Overview}
% ===========================================================================

The pion soft factor is a four-point correlation function that appears ubiquitously
in the factorization of transverse-momentum-dependent (qTMD) observables.
In the Collins--Soper (CS) framework~\cite{Collins:1981uk,Collins:1984kg} and the
Ji--Ma--Yuan (JMY) framework~\cite{Ji:2004wu,Ji:2019ewn}, the soft factor absorbs
rapidity divergences that arise when one attempts to separate the hard scattering,
the collinear beam functions, and the soft radiation in hadron scattering processes.

For a pion, the soft factor is defined through a matrix element of four Wilson-line
staple operators arranged on the light-cone.  In a Euclidean lattice calculation one
gives up the light-cone separation and instead works with equal-time correlators in
a large-momentum frame, following the LaMET/quasi-PDF philosophy.  The Lorentz
invariant definition is then replaced by a quasi-soft-factor whose infrared
structure coincides with the physical soft factor after a perturbative matching.

Concretely, the lattice object we compute is

\begin{equation}\label{eq:soft_def}
S(b_{\perp}, t; t_{0}, t_{\mathrm{sep}}, \Gamma_{1}, \Gamma_{2})
 \;=\;
\sum_{\mathbf{x}}
\Tr\Big[
  A_{b}(\mathbf{x}, t)\;\Gamma_{2}\;B_{b}(\mathbf{x}, t)\;\Gamma_{1}
\Big]\, ,
\end{equation}

where $b_{\perp}$ is a transverse displacement between the quark and antiquark
Wilson lines, $t_{0}$ is the source time slice, $t_{\mathrm{sep}}$ is the sink--source
separation, and the closed spin-color blocks $A_{b}$ and $B_{b}$ are built from four
Coulomb-gauge-fixed wall-source propagators as described in detail below.  The
matrices $\Gamma_{1}$ and $\Gamma_{2}$ are drawn from a reduced gamma set appropriate
for the soft factor.

The physical interpretation of Eq.~\eqref{eq:soft_def} is that the two blocks
$A_{b}$ and $B_{b}$ represent the quark and antiquark lines that travel from the
source to the sink, with the transverse displacement $b$ inserted to probe the
transverse-momentum dependence of the soft radiation.  The trace over spin and
color indices projects onto the desired gamma structures.

% ===========================================================================
\section{Two-Stage Workflow}
% ===========================================================================

The calculation is intentionally split into two independent stages to allow flexible
reuse of the expensive propagator inversions.

\paragraph{Stage 1 -- Propagator Generation}
(\walltag{Pyquda\_pion\_soft\_factor\_prop.py})

\begin{itemize}[nosep]
  \item Read a Coulomb-gauge-fixed configuration.
  \item For every requested source time slice $t_{0}$ and every quark momentum
        $\mathbf{k}$ (and $-\mathbf{k}$), construct a momentum-phase wall source and
        invert the Dirac operator.
  \item Save each propagator as an HDF5 file with full metadata (lattice tag,
        configuration number, smearing tag, source position, quark momentum).
\end{itemize}

\paragraph{Stage 2 -- Contraction}
(\walltag{Pyquda\_pion\_soft\_factor\_contract.py})

\begin{itemize}[nosep]
  \item Load the wall propagators saved in Stage~1 (no gauge field needed except for
        lattice geometry metadata).
  \item For each pair of forward/backward quark momenta $(\mathbf{k}_{\mathrm{fw}},
        \mathbf{k}_{\mathrm{bw}})$ compute:
        \begin{enumerate}[nosep,label=(\roman*)]
          \item the wall-to-wall two-point diagnostic $C_{2}(t)$,
          \item the qTMDWF-like diagnostic $C_{\mathrm{TMDWF}}(b_{T}, b_{z}, t)$,
          \item the full soft-factor four-point correlator
                $S(b_{\perp}, t; t_{\mathrm{sep}}, \Gamma_{1}, \Gamma_{2})$.
        \end{enumerate}
  \item Save all correlators to HDF5 with appropriate group hierarchies.
\end{itemize}

This two-stage design means that one can change contraction parameters
(e.g.\ $\Gamma_{1}$, $\Gamma_{2}$, $b_{T}$ ranges, sink interpolators) without
re-running any inversions.

% ===========================================================================
\section{Wall-Source Propagators}
% ===========================================================================

\subsection{Coulomb-Gauge Wall Source}

The source for quark momentum $\mathbf{k}$ at time $t_{0}$ is

\begin{equation}\label{eq:wall_src}
\eta_{\mathbf{k}}(\mathbf{x}, t)
   \;=\; \delta_{t,\,t_{0}}\;\exp\!\big( +i\,\mathbf{k}\!\cdot\!\mathbf{x} \big)\, .
\end{equation}

The spatial part is a plane wave with momentum $+\mathbf{k}$, applied uniformly on
every spatial site of time slice $t_{0}$.  The Coulomb gauge condition is imposed
before the wall source is constructed, which suppresses the gauge noise that would
otherwise contaminate a point source.  The corresponding propagator is

\begin{equation}\label{eq:wall_prop}
G_{\mathbf{k}}(x; t_{0}) \;=\; D^{-1}\,\eta_{\mathbf{k}}\, ,
\end{equation}

where $D$ is the Wilson-clover Dirac operator with appropriate boundary conditions
($t$-boundary set to antiperiodic, \walltag{t\_boundary = -1}).

\subsection{Both Signs of Quark Momentum}

The soft-factor contraction requires both $+\mathbf{k}$ and $-\mathbf{k}$
propagators on all time slices.  Stage~1 therefore loops over a set of momenta that
includes both signs:

\begin{equation}
\mathcal{K}_{\text{save}} \;=\;
  \{\,\mathbf{k}_{\mathrm{fw}},\,\mathbf{k}_{\mathrm{bw}},\,
   -\mathbf{k}_{\mathrm{fw}},\,-\mathbf{k}_{\mathrm{bw}}\,\}
\quad\text{(duplicates removed)}.
\end{equation}

\subsection{Antiquark Line from Gamma5 Hermiticity}

The Euclidean Dirac operator satisfies $\gamma_{5}$-hermiticity,

\begin{equation}\label{eq:gamma5_herm}
D^{\dagger} \;=\; \Gfive\, D\, \Gfive\, .
\end{equation}

For an antiquark propagating \emph{backward} in time the required propagator is
obtained from the $-\mathbf{k}$ forward propagator via

\begin{equation}\label{eq:antiquark}
\gbar{G}_{-\mathbf{k}}(x; t_{0})
   \;=\; \Gfive \; G_{-\mathbf{k}}(x; t_{0})^{\dagger} \; \Gfive\, .
\end{equation}

In the code this operation is performed by the \einsum\ pattern

\begin{verbatim}
    xp.einsum("ij, tzyxmlca, kl -> tzyxkjca",
             gamma5, prop.conj(), gamma5)
\end{verbatim}

which transposes the spin indices $(i\leftrightarrow j,\;k\leftrightarrow l)$
appropriately.  The color index $a$ and the spatial lattice indices $(z,y,x)$ are
spectators in this transformation.

\subsection{Momentum Phase Convention}

The phase convention is tied to the definition of the wall source, Eq.~\eqref{eq:wall_src}:

\begin{center}
\begin{tabular}{ll}
\toprule
Operation & Phase factor \\
\midrule
Wall source at $t_{0}$ with momentum $\mathbf{k}$
   & $\exp(+i\,\mathbf{k}\!\cdot\!\mathbf{x})$ \\[4pt]
Sink-side phase on $G_{\mathbf{k}}$ (two-point)
   & $\exp(-i\,\mathbf{p}\!\cdot\!\mathbf{x})$ where $\mathbf{p}$ is the meson momentum \\[4pt]
Sink-side phase on $G_{\mathbf{k}_{\mathrm{bw}}}$ (soft factor)
   & $\exp(-2i\,\mathbf{P}\!\cdot\!\mathbf{x})$ where $\mathbf{P}=\mathbf{k}_{\mathrm{fw}}+\mathbf{k}_{\mathrm{bw}}$ \\
\bottomrule
\end{tabular}
\end{center}

The sign of the applied phase in \walltag{apply\_phase} is controlled by the
\walltag{sign} argument.  When \walltag{sign=1}, the raw momentum-phase factor
$\exp(+i\mathbf{k}\!\cdot\!\mathbf{x})$ is multiplied onto the propagator data.
When \walltag{sign=-1}, the complex conjugate $\exp(-i\mathbf{k}\!\cdot\!\mathbf{x})$
is used instead.

% ===========================================================================
\section{Wall-to-Wall Two-Point Diagnostic}
% ===========================================================================

Before computing the four-point soft factor we verify that the wall-source
propagators produce a sensible pion two-point function.  Using the forward
propagator $G_{\mathrm{fw}} \equiv G_{\mathbf{k}}(t_{0})$ and the backward-source
propagator $G_{\mathrm{bw}} \equiv G_{-\mathbf{k}_{\mathrm{bw}}}(t_{0})$, the
wall-to-wall two-point function for source interpolator $\Gamma_{\mathrm{src}}$ and
sink interpolator $\Gamma_{\mathrm{sink}}$ is

\begin{equation}\label{eq:C2}
C_{2}(t) \;=\;
\sum_{\mathbf{x}} \Tr\Big[
   \Gamma_{\mathrm{src}}\;
   \gbar{G}_{\mathrm{bw}}(\mathbf{x}, t; t_{0})\;
   \Gamma_{\mathrm{sink}}\;
   G_{\mathrm{fw}}^{\mathrm{phased}}(\mathbf{x}, t; t_{0})
\Big]\, ,
\end{equation}

where

\begin{itemize}[nosep]
  \item $\gbar{G}_{\mathrm{bw}}$ is the gamma5-hermiticity-conjugated antiquark line,
        Eq.~\eqref{eq:antiquark}.
  \item $G_{\mathrm{fw}}^{\mathrm{phased}}$ is the forward propagator multiplied by
        the meson momentum phase $\exp(-i\mathbf{p}\!\cdot\!\mathbf{x})$ with
        $\mathbf{p} = \mathbf{k}_{\mathrm{fw}} + \mathbf{k}_{\mathrm{bw}}$.
  \item The backward propagator is \emph{not} phased, because the wall source
        $\eta_{-\mathbf{k}_{\mathrm{bw}}}(t_{0})$ already carries
        $-\mathbf{k}_{\mathrm{bw}}$ in its construction; the antiquark line thus
        naturally projects onto total momentum $\mathbf{p}$.
\end{itemize}

In the code, \walltag{contract\_wall\_2pt} performs this computation using
three nested \einsum\ calls:

\begin{enumerate}[nosep]
  \item $\gbar{G}_{\mathrm{bw}}$ via gamma5 conjugation.
  \item $\Gamma_{\mathrm{src}} \cdot \gbar{G}_{\mathrm{bw}} \cdot \Gamma_{\mathrm{sink}}$.
  \item Trace with $G_{\mathrm{fw}}^{\mathrm{phased}}$ and spatial sum.
\end{enumerate}

The final index structure is \walltag{(sink\_key, src\_key, t)}, and after gathering
across MPI ranks the time axis is rolled so that the source time $t_{0}$ maps to
$t=0$.

% ===========================================================================
\section{qTMDWF-like Diagnostic Check}
% ===========================================================================

The qTMDWF diagnostic generalizes the two-point function by shifting the backward
antiquark line transversely and longitudinally before contracting.  For transverse
displacement $b_{T}$ along axis $\hat{n}_{T}$ and longitudinal displacement $b_{z}$:

\begin{equation}\label{eq:tmdwf}
C_{\mathrm{TMDWF}}(b_{T}, b_{z}, t) \;=\;
\sum_{\mathbf{x}} \Tr\Big[
   \Gamma_{\mathrm{src}}\;
   \gbar{G}_{\mathrm{bw}}^{\,\mathrm{shift}}(\mathbf{x}, t; t_{0})\;
   G_{\mathrm{fw}}^{\mathrm{phased}}(\mathbf{x}, t; t_{0})
\Big]\, ,
\end{equation}

where $\gbar{G}_{\mathrm{bw}}^{\,\mathrm{shift}}$ is the gamma5-conjugated backward
propagator shifted by $(b_{T}, b_{z})$.  The forward propagator is left unshifted.

This diagnostic is valuable because:

\begin{enumerate}[nosep]
  \item It tests that the backward propagator carries the expected momentum
        and spatial structure.
  \item It previews the transverse displacement needed for the full soft factor.
  \item For $b_{T}=0,\;b_{z}=0$ it reduces to a simpler trace that can be compared
        against the two-point diagnostic as a consistency check.
\end{enumerate}

In the code, \walltag{contract\_tmdwf\_check} iterates over $b_{T}$ and $b_{z}$,
applies coordinate-gauge shifts via \walltag{prop.shift()}, gamma5-conjugates the
shifted propagator, and contracts against the phased forward propagator with a
single $\Gamma_{\mathrm{src}}$ insertion (no sink interpolator).

% ===========================================================================
\section{Soft-Factor Four-Point Contraction}
% ===========================================================================

This is the main measurement.  For each source time $t_{0}$ and sink separation
$t_{\mathrm{sep}}$ the contraction uses four Coulomb-gauge wall propagators,
organized by their role in the forward/backward quark and antiquark lines:

\begin{equation}\label{eq:four_props}
\begin{aligned}
G_{w}            &\equiv G_{+\mathbf{k}_{\mathrm{fw}}}(t_{0})
                  \qquad\text{(forward quark, source time)} \\[2pt]
G_{w,\mathrm{b}_{\perp}}^{\dagger}
                 &\equiv G_{-\mathbf{k}_{\mathrm{bw}}}(t_{0})
                  \qquad\text{(backward antiquark, source time)} \\[2pt]
G_{w}^{\dagger}  &\equiv G_{-\mathbf{k}_{\mathrm{fw}}}(t_{0}+t_{\mathrm{sep}})
                  \qquad\text{(forward antiquark, sink time)} \\[2pt]
G_{w}^{\mathrm{b}_{\perp}}
                 &\equiv G_{+\mathbf{k}_{\mathrm{bw}}}(t_{0}+t_{\mathrm{sep}})
                  \qquad\text{(backward quark, sink time)}
\end{aligned}
\end{equation}

The naming convention follows the legacy code:

\begin{itemize}[nosep]
  \item \walltag{Gw}: \textbf{G}auge-fixed \textbf{w}all propagator for the forward
        quark at the source.
  \item \walltag{Gw\_bperp\_dagger}: The dagger emphasizes that this propagator will
        be gamma5-conjugated (``dagger'') and that it corresponds to the backward
        line that will be transversely displaced (``bperp'').
  \item \walltag{Gw\_dagger}: Forward antiquark at the sink, gamma5-conjugated.
  \item \walltag{Gw\_bperp}: Backward quark at the sink, transversely displaced.
\end{itemize}

\subsection{Sink-Side Momentum Phase}

The backward quark propagator at the sink, $G_{w}^{\mathrm{b}_{\perp}}$, is multiplied
by the momentum-transfer phase

\begin{equation}\label{eq:sink_phase}
G_{w}^{\mathrm{b}_{\perp}}
\;\longrightarrow\;
\exp\!\big({-2i\,\mathbf{P}\!\cdot\!\mathbf{x}}\big)\;
G_{w}^{\mathrm{b}_{\perp}} ,
\qquad
\mathbf{P} \equiv \mathbf{k}_{\mathrm{fw}} + \mathbf{k}_{\mathrm{bw}}
\end{equation}

where $\mathbf{P}$ is the total pion momentum.  The factor of~2 arises from the
kinematic structure of the soft factor, where both the source and sink sides
contribute equal and opposite phases to the four-point function.  All other
propagators carry their momentum through the wall-source construction; only
$G_{w}^{\mathrm{b}_{\perp}}$ receives this explicit phase because it is the only
propagator for which the wall-source momentum ($+\mathbf{k}_{\mathrm{bw}}$) does not
match the required total momentum projection.

\subsection{Time-Major Lexicographic Layout}

Before contraction, all four propagators are rearranged into \emph{time-major}
lexicographic order via

\begin{verbatim}
    prop.lexico(False)
\end{verbatim}

This yields the index ordering

\begin{equation}
\underbrace{t}_{\text{time}}\,
\underbrace{z}_{\text{slowest spatial}}\,
\underbrace{y}_{\text{middle spatial}}\,
\underbrace{x}_{\text{fastest spatial}}\,
\underbrace{j,\,i}_{\text{spin sink, source}}\,
\underbrace{b,\,a}_{\text{color sink, source}}
\end{equation}

or, in the einsum convention used throughout the code,

\begin{equation}
\text{\walltag{tzyx\;j i\;b a}}
\end{equation}

where Latin letters $(i,j,k,l)$ denote Dirac spin indices and $(a,b,c)$ denote
color indices.  This time-major layout is essential because the final trace over
spatial volume and the time projection must respect the lexicographic order for
correct MPI halo exchanges.

% ===========================================================================
\section{Closed Spin-Color Blocks}
% ===========================================================================

For each transverse displacement vector $\mathbf{b}$ (pointing along the axis
$\hat{n}_{\mathrm{T}}$ with magnitude $b_{T}$ lattice units), we form two
closed spin-color blocks:

\subsection{Block $A_{b}$ -- Source Side}

\begin{equation}\label{eq:A_block}
A_{b}(\mathbf{x}, t) \;=\;
G_{w}(\mathbf{x}, t)\;
\cdot\;
\Gamma_{\mathrm{src}}\;
\cdot\;
\Gfive\;
\cdot\;
G_{w}^{\mathrm{b}_{\perp}\dagger}
   (\mathbf{x}+\mathbf{b},\,t)^{\dagger}
\cdot\;
\Gfive\, .
\end{equation}

\noindent
In words:
\begin{enumerate}[nosep]
  \item Start with the forward quark propagator $G_{w}$ at position $\mathbf{x}$.
  \item Insert the source interpolator $\Gamma_{\mathrm{src}}$.
  \item Gamma5-conjugate the transversely-shifted backward antiquark propagator
        $G_{w}^{\mathrm{b}_{\perp}\dagger}$ (the dagger in the name refers to the
        fact that this propagator is stored as $G_{-\mathbf{k}_{\mathrm{bw}}}$
        and must be conjugated).
  \item The shift by $+\mathbf{b}$ is applied to $G_{w}^{\mathrm{b}_{\perp}\dagger}$
        before conjugation.
\end{enumerate}

\subsection{Block $B_{b}$ -- Sink Side}

\begin{equation}\label{eq:B_block}
B_{b}(\mathbf{x}, t) \;=\;
G_{w}^{\mathrm{b}_{\perp}}(\mathbf{x}+\mathbf{b},\,t)
\cdot\;
\Gamma_{\mathrm{sink}}
\cdot\;
\Gfive\;
\cdot\;
G_{w}^{\dagger}(\mathbf{x},\,t)^{\dagger}
\cdot\;
\Gfive\, .
\end{equation}

\noindent
In words:
\begin{enumerate}[nosep]
  \item Start with the transversely-shifted backward quark propagator
        $G_{w}^{\mathrm{b}_{\perp}}$ at position $\mathbf{x}+\mathbf{b}$.
        This propagator already carries the sink-side phase,
        Eq.~\eqref{eq:sink_phase}.
  \item Insert the sink interpolator $\Gamma_{\mathrm{sink}}$.
  \item Gamma5-conjugate the forward antiquark propagator $G_{w}^{\dagger}$
        at position $\mathbf{x}$.
\end{enumerate}

Both blocks $A_{b}$ and $B_{b}$ are matrices in spin-color space at each spacetime
point $(\mathbf{x}, t)$.  Their open indices are the spin and color indices that
will be contracted in the final trace.

\subsection{Four-Point Correlator}

With the blocks defined, the soft-factor four-point function is assembled as

\begin{equation}\label{eq:C4_full}
\begin{aligned}
C_{4}(t; t_{\mathrm{sep}}, \mathbf{b}, \Gamma_{1}, \Gamma_{2})
\;=\;
\sum_{\mathbf{x}}
\Tr\Big[
   A_{b}(\mathbf{x}, t)\;
   \Gamma_{2}\;
   B_{b}(\mathbf{x}, t)\;
   \Gamma_{1}
\Big]\, .
\end{aligned}
\end{equation}

The trace is over the spin and color indices that connect $A_{b}$ and $B_{b}$; the
spatial sum $\sum_{\mathbf{x}}$ projects onto zero (or total) spatial momentum after
the phase factors have been accounted for in the propagator definitions.

% ===========================================================================
\section{Gamma and Interpolator Conventions}
% ===========================================================================

\subsection{Soft-Factor Gamma Set}

The soft factor uses a reduced set of gamma structures compared to standard meson
spectroscopy.  The PyQUDA gamma matrices follow the Euclidean conventions of the
underlying library:

\begin{center}
\begin{tabular}{lll}
\toprule
Key & Matrix & Description \\
\midrule
\walltag{"5"}  & $\Gfive = \gamma_{5}$          & Chirality \\
\walltag{"I"}  & $\gI$                          & Identity \\
\walltag{"X"}  & $\Gone$                        & $\gamma_{1}$ \\
\walltag{"Y"}  & $\Gtwo$                        & $\gamma_{2}$ \\
\walltag{"X5"} & $\Gone\Gfive$                  & $\gamma_{1}\gamma_{5}$ \\
\walltag{"Y5"} & $\Gtwo\Gfive$                  & $\gamma_{2}\gamma_{5}$ \\
\bottomrule
\end{tabular}
\end{center}

In the code these are stored in the dictionary \walltag{soft\_factor\_pyquda\_gammas},
which maps each key to the corresponding PyQUDA \walltag{gamma} object.  The
\walltag{gamma(15)} denotes $\Gfive$ (the 16th element in the standard Euclidean
gamma basis, 0-indexed), and \walltag{gamma(0)} is the identity.

\subsection{Pion Interpolators}

The default pion interpolator set is

\begin{equation}\label{eq:pion_interp}
\walltag{"Z5-X5"} \;=\; \Gthree\Gfive - \Gone\Gfive
   \;=\; \gamma_{3}\gamma_{5} - \gamma_{1}\gamma_{5}\, .
\end{equation}

This is a linear combination that projects onto the pseudoscalar pion with good
overlap at finite momentum.  Both $\Gamma_{1}/\Gamma_{2}$ (the soft-factor gamma
insertions) and \walltag{pion\_src}/\walltag{pion\_sink} (the pion interpolators)
are passed as Python dictionaries, so the user can supply arbitrary operator sets.

The legacy convention is that the source and sink interpolator dictionaries
are the same object by default, but they can be set independently to explore
operator dependence.

\subsection{Gamma5 as \texorpdfstring{$\Gfive$}{gamma5}}

The matrix $\Gfive$ appears in three distinct roles:

\begin{enumerate}[nosep]
  \item \textbf{Gamma5-hermiticity conjugation} of antiquark propagators,
        Eq.~\eqref{eq:antiquark}.
  \item \textbf{Block closure} in $A_{b}$ and $B_{b}$, Eqs.~\eqref{eq:A_block}--\eqref{eq:B_block}.
  \item \textbf{As a gamma insertion}: when $\Gamma_{1}$ or $\Gamma_{2}$ is
        \walltag{"5"}, it acts as a genuine operator insertion probing the chiral
        structure of the soft factor.
\end{enumerate}

In the code, the variable \walltag{G5 = gamma.gamma(15)} is used for roles~1 and~2,
while role~3 uses the entry \walltag{soft\_factor\_pyquda\_gammas["5"]}, which is
the same underlying matrix.

% ===========================================================================
\section{Transverse Displacement Implementation}
% ===========================================================================

\subsection{Coordinate-Gauge Shifts}

The transverse displacement $\mathbf{b}$ is implemented by a simple array roll
in the coordinate gauge (CG).  No explicit gauge link is inserted between the
shifted quark fields.  This is consistent with the CG convention used in the
pion qTMD workflow: in Coulomb gauge the gluon field is minimized, and at
leading order the gauge link reduces to unity.

For a displacement of $b_{T}$ lattice units along axis $\hat{n}_{T}$ (where
$n_{T}=0$ means $x$, $n_{T}=1$ means $y$, $n_{T}=2$ means $z$), the operation is

\begin{verbatim}
    xp.roll(Gw_bperp, shift=bT, axis=axis)
\end{verbatim}

where \walltag{axis = 3 - bT\_dir} converts from the physics axis convention
to the time-major lexicographic index convention.  In the time-major layout
\walltag{tzyx...}, the spatial axes are ordered $(z=1, y=2, x=3)$, so:

\begin{center}
\begin{tabular}{ll}
\toprule
\walltag{bT\_dir} & \walltag{axis} in lexicographic layout \\
\midrule
0 ($x$) & 3 \\
1 ($y$) & 2 \\
2 ($z$) & 1 \\
\bottomrule
\end{tabular}
\end{center}

\subsection{Displacement Parameters}

\begin{itemize}[nosep]
  \item \walltag{bT\_dir}: list of axes along which to displace, e.g.\
        \walltag{[0]} for $x$ only, \walltag{[0, 1]} for both $x$ and $y$.
  \item \walltag{bT\_length}: maximum displacement in lattice units (inclusive;
        loop runs from $0$ to \walltag{bT\_length}).
  \item \walltag{bz\_length}: maximum longitudinal displacement (default $0$,
        used only in the qTMDWF diagnostic).
\end{itemize}

% ===========================================================================
\section{Code Implementation}
% ===========================================================================

The core logic lives in the class \walltag{pion\_soft\_factor} (file
\walltag{pion\_soft\_factor\_vibe\_develop.py}).  This section documents each method.

\subsection{\walltag{\_\_init\_\_(parameters)}}

Accepts a dictionary with keys:

\begin{center}
\begin{tabular}{ll}
\toprule
Key & Description \\
\midrule
\walltag{quark\_mom} & List of quark 3-momenta $[\mathbf{k}_{1}, \mathbf{k}_{2}, \dots]$ \\
\walltag{bT\_dir}    & List of transverse displacement axes \\
\walltag{bT\_length} & Maximum $b_{T}$ \\
\walltag{bz\_length} & Maximum $b_{z}$ (qTMDWF diagnostic) \\
\walltag{tsep\_list} & List of sink separations \\
\walltag{pion\_channel\_pairs} & Ordered pair label $\mapsto(\Gamma_{\rm src},\Gamma_{\rm sink})$ \\
\walltag{gamma\_channel\_pairs} & Ordered pair label $\mapsto(\Gamma_1\text{ label},\Gamma_2\text{ label})$ \\
\bottomrule
\end{tabular}
\end{center}

The channel axes are paired rather than Cartesian products.  The stored order
is
\texttt{[pion\_pair,gamma\_pair,bT\_direction,bT,time]} (with an outer
\texttt{tsep} axis in the combined output).  Each pair label and both members
are written explicitly to HDF5.  The insertion labels are resolved through the
canonical raw PyQUDA basis
\texttt{5,I,X,Y,X5,Y5}; in particular \texttt{Y5} is the raw bit-mask matrix,
not a separately sign-adjusted physical axial matrix.  The file stores the raw
basis, PyQUDA IDs, and \texttt{physical\_from\_pyquda} transform.

\subsection{\walltag{create\_wall\_src(latt\_info, tslice, momentum)}}

Constructs the Coulomb-gauge wall source, Eq.~\eqref{eq:wall_src}.  Internally it
calls \walltag{phase.MomentumPhase} to build the momentum phase factor
$\exp(+i\mathbf{k}\!\cdot\!\mathbf{x})$ and \walltag{source.propagator} to place it
on the specified time slice.

\subsection{\walltag{create\_wall\_propagator(dirac, latt\_info, tslice, momentum)}}

Assembles the wall source and inverts the Dirac operator:
\walltag{core.invertPropagator(dirac, wall\_src, 1, 0)}.  The arguments \walltag{1, 0}
specify the multigrid level and the number of deflation vectors, respectively
(standard single-level inverter with no deflation).

\subsection{\walltag{save\_wall\_propagator(prop, tag, attrs=None)}}

Saves the propagator to HDF5 using \walltag{prop.saveH5(tag + ".h5", "propagator")}.
Optional metadata (lattice tag, configuration number, smearing, source time, quark
momentum) is written as HDF5 attributes for reproducibility.

\subsection{\walltag{load\_wall\_propagator(latt\_info, tag)}}

Loads a previously saved propagator via
\walltag{core.LatticePropagator.loadH5(tag + ".h5", "propagator")}.

\subsection{\walltag{apply\_phase(prop, momentum, sign=1, x0=None)}}

Multiplies the propagator data by a momentum phase factor:

\begin{itemize}[nosep]
  \item Computes the phase $\exp(\pm i\mathbf{k}\!\cdot\!(\mathbf{x}-\mathbf{x}_{0}))$
        using \walltag{MomentumPhase.getPhase}.
  \item \walltag{sign=1} gives the positive phase; \walltag{sign=-1} gives the
        complex conjugate.
  \item The phase array is broadcast over the spin-color trailing dimensions.
  \item Returns a copy of the propagator with the phased data.
\end{itemize}

\subsection{\walltag{contract\_wall\_2pt(latt\_info, prop\_fw, prop\_bw, pion\_mom, src\_key, sink\_keys)}}

Computes the wall-to-wall two-point function, Eq.~\eqref{eq:C2}.  Steps:

\begin{enumerate}[nosep]
  \item Phase the forward propagator with $\exp(-i\mathbf{p}\!\cdot\!\mathbf{x})$.
  \item Lexicographically reorder both propagators: \walltag{.lexico(False)}.
  \item Gamma5-conjugate the backward propagator.
  \item Contract with source and sink interpolators.
  \item Trace and spatial sum.
  \item Gather across MPI ranks.
\end{enumerate}

\subsection{\walltag{contract\_tmdwf\_check(latt\_info, prop\_fw, prop\_bw, pion\_mom, src\_key)}}

Computes the qTMDWF-like diagnostic, Eq.~\eqref{eq:tmdwf}.  Iterates over
$(b_{T}, b_{z})$ and returns a flat list of time-series correlators.

\subsection{\walltag{contract\_soft\_factor(latt\_info, prop\_fw, prop\_bw\_src, prop\_sink\_bw, prop\_sink\_fw, pion\_mom)}}

The main four-point contraction, described in full in Sec.~\ref{sec:einsum}.
Returns the five-dimensional array \walltag{(n\_src, n\_gm1, n\_dir, n\_bT+1, Lt)}
together with the ordered key lists.

% ===========================================================================
\section{Einsum Contraction Patterns}
\label{sec:einsum}
% ===========================================================================

The soft-factor contraction uses four \einsum\ calls per $(b_{T}, \Gamma_{1})$
combination.  This is the most complex contraction in the Pyquda\_Measurement
codebase.  We document the index conventions here.

\subsection{Index Convention}

In the time-major lexicographic layout, propagators carry the indices

\begin{equation}\label{eq:indices}
\underbrace{t}_{\text{time}}\,
\underbrace{z}_{\text{spatial }2}\,
\underbrace{y}_{\text{spatial }1}\,
\underbrace{x}_{\text{spatial }0}\,
\underbrace{j}_{\text{spin sink}}\,
\underbrace{i}_{\text{spin source}}\,
\underbrace{b}_{\text{color sink}}\,
\underbrace{a}_{\text{color source}}
\end{equation}

Gamma matrices are $4\times4$ with indices $(k,l)$ in spin space.  The
interpolator stacks $\Gamma_{\mathrm{src}}$ and $\Gamma_{\mathrm{sink}}$ carry
an additional leading batch index $s$ (enumerating the source or sink keys).
Similarly, the $\Gamma_{1}$ and $\Gamma_{2}$ stacks carry a batch index over
the gamma keys.

\subsection{Einsum 1: Closed Block \texorpdfstring{$A_{b}$}{A\_b}}

\begin{equation}\label{eq:einsum_A}
\begin{aligned}
A_{b} \;=\; &
\einsum\big(
   \text{\walltag{tzyxjiba, sik, kl, tzyxmlca, mn}}
   \to \text{\walltag{stzyxjnbc}}, \\
   &\qquad G_{w},\;
        \Gamma_{\mathrm{src}}^{\mathrm{stack}},\;
        \Gfive,\;
        \big[G_{w}^{\mathrm{b}_{\perp}\dagger}\big]^{\mathrm{shift},\,\conj{}},\;
        \Gfive
\big)
\end{aligned}
\end{equation}

Index mapping:

\begin{center}
\begin{tabular}{ll}
\toprule
Index & Carried by \\
\midrule
$t,z,y,x$ & Spacetime (summed/kept depending on position) \\
$j,i$     & Spin sink/source of $G_{w}$ \\
$b,a$     & Color sink/source of $G_{w}$ \\
$s$       & Batch over source interpolator keys \\
$k,l,m,n$ & Spin contraction indices (closed in the einsum) \\
$l,c$     & Spin/color of $G_{w}^{\mathrm{b}_{\perp}\dagger}$ (sink side) \\
$n,c$     & Output spin/color indices \\
\bottomrule
\end{tabular}
\end{center}

The superscript ``shift, $\conj{}$'' means that $G_{w}^{\mathrm{b}_{\perp}\dagger}$
has been rolled by $b_{T}$ along the appropriate spatial axis and then complex
conjugated (the dagger in the name).  The two $\Gfive$ insertions perform the
gamma5-hermiticity conjugation of the shifted propagator.

The output block $A_{b}$ has shape $(s, t, z, y, x, j, n, b, c)$, where $s$
indexes the source interpolator, $(j,n)$ are the open spin indices, and $(b,c)$
are the open color indices that will connect to $B_{b}$.

\subsection{Einsum 2: Closed Block \texorpdfstring{$B_{b}$}{B\_b}}

\begin{equation}\label{eq:einsum_B}
\begin{aligned}
B_{b} \;=\; &
\einsum\big(
   \text{\walltag{tzyxjiba, sik, kl, tzyxmlca, mn}}
   \to \text{\walltag{stzyxjnbc}}, \\
   &\qquad G_{w}^{\mathrm{b}_{\perp},\,\mathrm{shift}},\;
        \Gamma_{\mathrm{sink}}^{\mathrm{stack}},\;
        \Gfive,\;
        \big[G_{w}^{\dagger}\big]^{\conj{}},\;
        \Gfive
\big)
\end{aligned}
\end{equation}

This is structurally identical to $A_{b}$ but with different propagators and
interpolators:

\begin{itemize}[nosep]
  \item $G_{w}^{\mathrm{b}_{\perp},\,\mathrm{shift}}$ is the transversely shifted
        \emph{and sink-phased} backward quark propagator.
  \item $G_{w}^{\dagger}$ is the forward antiquark propagator from the sink
        time slice, complex conjugated but \emph{not} shifted.
  \item $\Gamma_{\mathrm{sink}}$ replaces $\Gamma_{\mathrm{src}}$.
\end{itemize}

\subsection{Einsum 3: Four-Point Trace}

\begin{equation}\label{eq:einsum_C4}
\begin{aligned}
C_{4}(\mathbf{x}, t) \;=\; &
\einsum\big(
   \text{\walltag{tzyxjiba, ik, tzyxklba, lj}}
   \to \text{\walltag{tzyx}}, \\
   &\qquad A_{b}[\text{isrc}],\;
        \Gamma_{2}[\text{igm}],\;
        B_{b}[\text{isrc}],\;
        \Gamma_{1}[\text{igm}]
\big)
\end{aligned}
\end{equation}

This einsum closes the spin and color indices:

\begin{itemize}[nosep]
  \item $j$ (spin of $A_{b}$)$\;\to\;$ $i$ (row of $\Gamma_{2}$)
  \item $k$ (column of $\Gamma_{2}$) $\;\to\;$ $k$ (spin of $B_{b}$)
  \item $l$ (spin of $B_{b}$) $\;\to\;$ $l$ (row of $\Gamma_{1}$)
  \item $j$ (column of $\Gamma_{1}$) $\;\to\;$ $j$ (spin of $A_{b}$)
  \item $b,a$ color indices are traced within \einsum.
\end{itemize}

The result is a scalar on spin-color space at each spacetime point $(\mathbf{x},t)$.

\subsection{Einsum 4: Time Projection}

\begin{equation}\label{eq:einsum_time}
C_{4}(t) \;=\;
\einsum\big(
   \text{\walltag{tzyx $\to$ t}},\;
   C_{4}(\mathbf{x}, t)
\big)
\end{equation}

This sums over the spatial volume $\mathbf{x} = (z,y,x)$ at each time slice,
yielding the time-series correlator $C_{4}(t)$.  After this step the result is
gathered across MPI ranks.

\subsection{Full Loop Structure}

The contraction loops are nested as follows (outermost first):

\begin{enumerate}[nosep]
  \item $b_{T}$ direction (\walltag{idir})
  \item $b_{T}$ magnitude (0 to \walltag{bT\_length})
  \item Source interpolator key (\walltag{isrc})
  \item Gamma-1 key (\walltag{igm})
\end{enumerate}

For each $(b_{T}, \text{dir})$ pair, the shifted propagators and the two blocks
$A_{b}$, $B_{b}$ are computed \emph{outside} the $(\text{isrc}, \text{igm})$ loops.
Only the final trace (Einsum~3) and time projection (Einsum~4) are inside the
inner loops.  This is a significant optimization because the block construction
dominates the computational cost.

% ===========================================================================
\section{HDF5 Output Structure}
% ===========================================================================

\subsection{Propagator Files}

\begin{verbatim}
data/pion_soft_factor_prop/
  {lat}.pion_soft_factor_prop.{cfg}.{ama}.{src}.{sm}.{mom}.h5
\end{verbatim}

Each file contains a single dataset \walltag{"propagator"} and HDF5 attributes:
\walltag{lat\_tag}, \walltag{config\_num}, \walltag{sm\_tag}, \walltag{tslice},
\walltag{quark\_momentum} (as int32 array).

\subsection{Two-Point Diagnostic}

\begin{verbatim}
data/pion_soft_factor_c2pt/
  {lat}.pion_soft_factor_c2pt.{cfg}.{ama}.{src}.{sm}.fw_{mom1}.bw_{mom2}.h5
\end{verbatim}

HDF5 group hierarchy (from \walltag{save\_pion\_soft\_factor\_c2pt\_hdf5\_noRoll}):

\begin{verbatim}
  /SS/{src_key}/{sink_key}/PX{p[0]}PY{p[1]}PZ{p[2]}   -> dataset
\end{verbatim}

\subsection{qTMDWF Diagnostic}

\begin{verbatim}
data/pion_soft_factor_qTMDWF/
  {lat}.pion_soft_factor_qTMDWF.{cfg}.{ama}.{src}.{sm}.fw_{mom1}.bw_{mom2}.h5
\end{verbatim}

HDF5 group hierarchy (from \walltag{save\_pion\_soft\_factor\_qTMDWF\_hdf5\_noRoll}):

\begin{verbatim}
  /SP/{src_key}/PX{px}PY{py}PZ{pz}/{b_dir}/bT{bT}/bz{bz}   -> dataset
\end{verbatim}

\subsection{Four-Point Soft Factor}

\begin{verbatim}
data/pion_soft_factor/
  {lat}.pion_soft_factor.{cfg}.{ama}.{src}.{sm}.fw_{mom1}.bw_{mom2}.h5
\end{verbatim}

HDF5 group hierarchy (from \walltag{save\_pion\_soft\_factor\_hdf5\_noRoll}):

\begin{verbatim}
  /src{src_key}_sink{sink_key}/{gm1}_{gm2}/{b_dir}_{bT}/ts{tsep}   -> dataset
\end{verbatim}

Each dataset contains a 1D complex array of length $L_{t}$, the time-series
correlator for that combination of source interpolator, gamma insertions,
transverse displacement, and sink separation.

% ===========================================================================
\section{Comparison with Legacy Code}
% ===========================================================================

This implementation is a port from the GPT/PyQUDA mixed workflow
\walltag{PyQUDA\_qTMD\_ff\_4pt\_einsum.py}.  The key differences are:

\begin{enumerate}[nosep]
  \item \textbf{Pure PyQUDA}: No GPT dependency.  All gauge fixing, smearing,
        inversion, and I/O use the PyQUDA stack.
  \item \textbf{HDF5 propagator storage}: Propagators are saved between Stage~1
        and Stage~2, enabling (a) decoupled production and analysis, (b) reuse
        of propagators for multiple contraction parameter sets, and
        (c) portability across machines (propagators can be analyzed locally).
  \item \textbf{Identical physics}: The contraction patterns, gamma conventions,
        phase factors, and wall-source definitions are preserved from the legacy
        workflow to ensure output compatibility.  The numerical results should
        agree to machine precision with the GPT/PyQUDA mixed code.
  \item \textbf{Class-based design}: The \walltag{pion\_soft\_factor} class
        encapsulates all parameters and provides methods for each computational
        step, improving maintainability and enabling unit testing.
  \item \textbf{Backend abstraction}: All \einsum\ calls are dispatched through
        \walltag{xp} (which can be \walltag{numpy}, \walltag{cupy}, \walltag{dpnp},
        or \walltag{torch}), so the same code runs on CPU, NVIDIA GPU, Intel GPU,
        or PyTorch backends without modification.  Gamma stacks are first moved
        to the backend and, for \walltag{dpnp}, the exact SYCL queue of the
        reference propagator before stacking.  Gathered arrays use the shared
        backend-aware host conversion rather than backend-specific
        \walltag{asnumpy} calls.
\end{enumerate}

% ===========================================================================
\section*{Acknowledgments}

This work builds on the PyQUDA framework and the legacy GPT/PyQUDA mixed
soft-factor code.  The implementation follows the conventions established
in the pion qTMD and qTMDWF analyses.

% ===========================================================================
\begin{thebibliography}{99}

\bibitem{Collins:1981uk}
J.~C.~Collins and D.~E.~Soper,
``Back-To-Back Jets in QCD,''
Nucl.\ Phys.\ B \textbf{193}, 381 (1981).

\bibitem{Collins:1984kg}
J.~C.~Collins and D.~E.~Soper,
``The Theorems of Perturbative QCD,''
Ann.\ Rev.\ Nucl.\ Part.\ Sci.\ \textbf{37}, 383 (1987).

\bibitem{Ji:2004wu}
X.~Ji, J.~Ma and F.~Yuan,
``QCD factorization for semi-inclusive deep-inelastic scattering at low
transverse momentum,''
Phys.\ Rev.\ D \textbf{71}, 034005 (2005)
[arXiv:hep-ph/0404183].

\bibitem{Ji:2019ewn}
X.~Ji, Y.~Liu and Y.-S.~Liu,
``Transverse-momentum-dependent parton distribution functions from
large-momentum effective theory,''
Phys.\ Lett.\ B \textbf{811}, 135946 (2020)
[arXiv:1911.03840 [hep-ph]].

\end{thebibliography}

\end{document}
